\clearpage
\setcounter{page}{1}
\maketitlesupplementary


\section{Detailed Literature Review}
\label{app:litreview}

This appendix expands the compressed related-work discussion of
Section~\ref{sec:related} into the five bodies of work the paper builds
on or departs from.

\subsection{Procedural plant models and inverse procedural modeling}
Procedural modeling of plants begins with
L-systems~\citep{prusin1990abop} and matured into generators whose
parameters are organ-level, biologically meaningful
quantities~\citep{floraforge2026,hadadi2025procedural}. Inverting such
models, that is, recovering their parameters from observations, has a
long history for woody species: Stava et al.\ fit tree parameters by stochastic
search~\citep{stava2014ipmtrees}, Talton et al.\ formulated procedural
inversion as MCMC inference~\citep{talton2011metropolis}, and Guo et al.\
inferred L-systems directly from point clouds~\citep{guo2020ipmlsys},
with neural variants following~\citep{latentlsys2023}. These methods
optimize or infer per instance; amortized, feed-forward inversion of a
plant generator from field-degraded input is the gap this paper fills.
The closest amortized precedents sit outside agriculture: shape-program
induction~\citep{shapeprograms2019}, diffusion-based inverse procedural
content generation from a single image~\citep{zhao2025dipcg}, and
cadrille, which maps images, point clouds, or text to executable CAD
programs~\citep{kolodiazhnyi2025cadrille,text2cad2024}. All operate on
rigid, isolated, clean-input objects; none confronts deformable slender
organs, inter-plant occlusion, or real-sensor noise, and none studies
whether the program parameters should participate in the encoder.

\subsection{Learned parametric recovery of crop plants}
Two 2025 systems recover parametric plants directly from scans. Demeter
learns a morphology model (topology, articulation, PCA shape) from real
soybean point clouds~\citep{demeter2025}; Ghrer et al.\ infer
grammar-style parameterized plants from 3D scans~\citep{ghrer2025paramplants}.
Both take a single, reasonably complete point cloud as input and neither
addresses missing modalities or row occlusion. CropCraft recovers crop
structure by Bayesian optimization of a procedural model against
multi-view NeRF depth~\citep{cropcraft2026}; the approach is accurate
but per-scene and slow. A staged pipeline annotates terrestrial
laser scans with a vision-language model and fits a NURBS procedural
descriptor by differentiable refinement~\citep{hadadi2026vlmmaize,
hadadi2025procmaize,nurbsdiff2022}; it handles isolated, occlusion-free
plants at ${\sim}10$ minutes each and serves in this paper as both the
occlusion-free baseline and the source of real-data procedural ground
truth. Tree-D Fusion reconstructs simulation-ready urban trees from
street imagery via diffusion priors~\citep{lee2024treedfusion}, but its
envelope-then-branch pipeline targets woody silhouettes, a different
geometry class from thin cereal blades.

\subsection{Point-cloud completion and slender structures}
Learned completion runs from PCN~\citep{yuan2018pcn} through
transformer approaches~\citep{yu2021pointr,yu2023adapointr} to recent
extrapolation--interpolation designs~\citep{cai2024einet}, trained and
evaluated on benchmarks of compact rigid objects. Neither
PCN's eight categories nor ShapeNet-55/34 contains a
slender-under-occlusion class, and completion systems assume the missing
geometry is self-occluded surface, absent competing clutter from
\emph{other} objects inside the crop region. Plant-specific completion
is emerging, with amodal segmentation plus completion for occluded
plants~\citep{lee2026occ}, but it outputs geometry rather than
structure or parameters. Curve-aware and thin-structure reconstruction
\citep{curvegaussian2025,qi2026rethinking} shows growing awareness of the
slender regime, on which the standard metric itself fails: Chamfer
distance under-weights thin sheets~\citep{wu2021dcd}, a failure now
documented analytically~\citep{song2026structural} and addressed by density-aware
and information-theoretic variants~\citep{wu2021dcd,lin2023infocd}. We
take the further step of pairing geometric metrics against frozen-latent
phenotype probes and show they can rank modalities oppositely.

\subsection{3D plant phenotyping}
Organ-level 3D phenotyping datasets, including Pheno4D's labeled maize
and tomato sequences~\citep{schunck2021pheno4d},
Crops3D~\citep{crops3d2024}, and field-scale maize scans with procedural
fits~\citep{maizefield3d2025}, anchor trait extraction from point
clouds~\citep{qi2017pointnet,qi2017pointnetpp}. Radiance-field capture is
increasingly practical in field conditions~\citep{arshad2024nerf_field}.
No existing dataset, however, combines real field capture, slender
organs, inter-plant row occlusion, and procedural ground truth; our
benchmark is designed to close exactly that conjunction.

\subsection{Multimodal masked pretraining and distillation}
Masked autoencoding~\citep{he2022mae} extends naturally to multiple
modalities: MultiMAE over RGB, depth, and semantics~\citep{bachmann2022multimae};
4M and 4M-21 to any-to-any generation over tokenized
modalities~\citep{mizrahi2023fourm,bachmann2024fourm21}; Point-MAE to
geometry~\citep{pang2022pointmae}; and EmbodiedMAE, our architectural
ancestor, to RGB, depth, and point clouds for robot
manipulation~\citep{dong2025embodiedmae}. Missing-modality robustness is
an active concern~\citep{wang2024robustmissing}, and cross-modal
distillation from multi-sensor teachers to single-sensor students is
established for discriminative tasks; we distill a \emph{generative}
full-modality teacher for structured-parameter regression, and encode
continuous parameters with xVal-style number
tokens~\citep{golkar2023xval}. Our finding that a privileged parameter
stream suppresses representation quality connects to modality
competition~\citep{peng2022ogm} and gradient
conflict~\citep{yu2020pcgrad,chen2018gradnorm}, which supply candidate
mechanisms; no prior work, to our knowledge, documents the effect for
generator parameters in multimodal pretraining.

\subsection{Synthetic-data pretraining and sim-to-real}
Pretraining without natural data is viable, from abstract
fractals~\citep{kataoka2020fractaldb} to procedural
worlds~\citep{raistrick2023infinigen,greff2022kubric}, and domain
randomization bridges the reality gap~\citep{tobin2017domainrand}, with
synthetic value known to be task-dependent~\citep{fan2023intriguing} and
self-consumption a documented risk~\citep{alemohammad2023mad}. Our
setting is the favorable case: the deployment target \emph{is} the
generator's parameterization, the generator spans the morphology of the
crop, and real-field evaluation closes the loop.

% =========================================================================
\section{Detailed Ablation Studies}
\label{app:ablations}

All numbers in this appendix are from the 15K-plant studies
(\oldnum{purple} throughout); every table will be refreshed from the
10K--100K reruns. Protocol: every arm trains at a matched budget of
\oldnum{197{,}400} optimizer steps, global batch 32; representation
quality is a ridge probe on the frozen CLS token, one fixed view per
plant, identical target vectors across arms; generation quality is
validation Chamfer / parameter MAE at the matched checkpoint.

\subsection{Modality value-add (E2)}
Four pretraining arms differ only in which token streams exist.
\begin{table}[h]\centering\small
\begin{tabular}{lcc}
\toprule
Arm & PC Chamfer $\downarrow$ & vs.\ PC-only \\
\midrule
PC only                  & \oldnum{0.00495}  & ---            \\
PC + RGB                 & \oldnum{0.00250}  & \oldnum{$-49.5\%$} \\
PC + RGB + depth         & \oldnum{0.00171}  & \oldnum{$-65.6\%$} \\
PC + RGB + depth + params& \oldnum{0.00248}  & \oldnum{$-50.0\%$} \\
\bottomrule
\end{tabular}
\caption{Sorghum E2: adding the parameter stream to the encoder returns
Chamfer to the PC+RGB level, surrendering the depth gain. The frozen
probe agrees: the four-modality arm is worse on all eight sorghum
phenotype targets and \oldnum{7 of 11} maize targets.
\todo{add maize table + probe $R^2$ columns + CIs}}
\label{tab:e2}
\end{table}

\subsection{Token-budget control arms}
The four-stream arm shows each vision stream fewer visible tokens, a
mechanical confound. Control arms restore RGB/depth/PC to exactly the
three-modality arm's \oldnum{117} visible vision tokens, varying only
parameter-text visibility. With the budget restored the gap persists
(height probe \oldnum{0.958} vs.\ \oldnum{0.981}; roll \oldnum{0.764}
vs.\ \oldnum{0.806}), and visibility is dose-dependent: at 60\% text
visible, every skilled target worsens further. The effect is real, and
the budget explanation is excluded. \todo{full control-arm table with
bootstrap CIs}

\subsection{Parameters as a target only}
Masking parameter text at 100\% (never an input; reconstructed only by
the decoder) still loses to omitting the stream entirely: size targets
\oldnum{$-0.010$ to $-0.012$} $R^2$, roll \oldnum{$-0.069/-0.077$}
(val/test), mean leaf length \oldnum{$-0.102/-0.097$}, CIs clear of zero.
The latent analysis locates the damage: the target-only arm's latent is
the most size-dominated of its cohort (PC1 $=$ size at $|r|{=}$
\oldnum{0.95}, effective rank \oldnum{4.6}). The reconstruction pressure
itself, under the encoder's capacity, is what costs structure.
\todo{reconcile with matched-LR pc2text control when it lands}

\subsection{Data and model scaling at fixed compute}
\begin{table}[h]\centering\small
\begin{tabular}{lcccc}
\toprule
 & \multicolumn{2}{c}{Sorghum} & \multicolumn{2}{c}{Maize} \\
Train plants & Chamfer & Param acc@.05 & Chamfer & Probe moves? \\
\midrule
1{,}000  & \oldnum{0.00299} & \oldnum{0.608} & \oldnum{0.00292$^{*}$} & --- \\
3{,}000  & \oldnum{0.00258} & \oldnum{0.775} & \oldnum{$-34\%$} & yes \\
10{,}000 & \oldnum{0.00241} & \oldnum{0.811} & \oldnum{$-34\%$} & $\le$ noise \\
\bottomrule
\end{tabular}
\caption{Fixed 197{,}400-step budget. Near-duplicate 10{,}000 vs.\
10{,}500 arms give the run-to-run noise floor (\oldnum{$\sim$3\%}
Chamfer, \oldnum{0.02} $R^2$). Chamfer keeps paying for data; the probe
saturates by $\sim$3{,}000 plants in both species. More data buys
geometry, and stops buying linearly decodable phenotype.
$^{*}$maize small-data arms actively overfit (best at 29\% of schedule,
+79\% by the end). \todo{refresh at 100K; add model-scaling row (25M
$\to$ 114M $=-17.1\%$; 114M $\to$ 332M $\approx$ 0: base is the knee)}}
\label{tab:scaling}
\end{table}

\subsection{Distillation, per source modality}
Distilling the full-modal teacher into the warm-started student improves
the mean generation metric \oldnum{19.1\%}, concentrated where the
student was weakest: conditioning on PC alone \oldnum{$-27.5\%$} (its
Chamfer \oldnum{$-56.6\%$}), on parameters alone Chamfer
\oldnum{$-61.9\%$}; RGB and depth, already strong, gain
\oldnum{$\sim$11\%}. PC$\rightarrow$params MAE falls \oldnum{43.8\%}
(\oldnum{0.0617} $\to$ \oldnum{0.0347}), and distilling all four targets
beats distilling the parameter head alone (\oldnum{$-43.8\%$} vs.\
\oldnum{$-36.9\%$}); the other heads act as auxiliary supervision.
Distillation moves the whole view-difficulty curve down without
flattening it (\oldnum{9/10} views improve; best-to-worst ratio
unchanged), so view robustness needs a different lever.
\todo{LR caveat: specialist ran at 3e-5 vs.\ 1e-4; matched-LR control
pending. Per-trait recovery tables (PC vs.\ RGB conditioning,
position-only baseline) to follow from reruns.}

\subsection{View invariance without labels}
Ten views of one plant share one parameter ground truth, so the spread
of ten predictions is a label-free view-invariance measure. Distillation
roughly halves error \emph{and} view spread on the fields with skill
(starting point \oldnum{$-59\%$}, branching angle \oldnum{$-60\%$});
zero-skill fields instead grow \emph{more} view-dependent; the head
starts tracking image noise. ``Stable'' and ``right'' are separate
axes: several fields move less than 0.3 of their spread across views
while remaining worse than predicting the plant mean.

\subsection{Frozen-encoder baselines and camera-pose forensics}
Against four public frozen encoders under one probe
(EmbodiedMAE-B~\citep{dong2025embodiedmae},
Point-MAE~\citep{pang2022pointmae}, DINOv2~\citep{oquab2024dinov2},
MultiMAE~\citep{bachmann2022multimae}) plus a supervised-from-scratch
network, our three-modality arm wins every learnable sorghum target
(roll +\oldnum{0.078}, mean leaf length +\oldnum{0.104} over the
strongest baseline) and the margin comes from fusing RGB with geometry:
off-the-shelf Point-MAE beats our PC-only arm, and the four-modality
reference arm falls below EmbodiedMAE-B on size targets; this is the
parameter-input effect seen from outside. Supplying Point-MAE with each
view's full camera pose hands it the generator's fixed plant azimuth, a
canonicalization no field rig knows; restricted to gravity alone (what a
calibrated rig knows), our pose-free model wins \oldnum{5} of 9 decided
targets. Tassel droop stays below zero for every encoder of this design and
remains an open limitation. \todo{full tables, both species, with CIs}

\subsection{Latent structure (E9)}
The latents are continua, as they should be for a single-generator
population (silhouette $\le$ \oldnum{0.24}; k-means at chance against
leaf-count deciles); the question is factor alignment. Sorghum latents
occupy \oldnum{4--7} effective dimensions, maize \oldnum{7--17},
consistent with maize's more independent generative factors. The
dominant principal component is \emph{never} the phenotype
(\oldnum{$\sim$30\%} of variance unexplained by any trait), and linear
decodability dissociates from local embedding structure: roll angle reads
at $R^2$ \oldnum{0.806} from a latent where it is barely locally
organized. t-SNE smoothness is the wrong instrument for representation
claims on continuous generators; probes are the right one.
\todo{E9 figures; PC1 identity remains open}